% TOP_PROBLEM: {"id": 3242, "title": "Melnikov's valency-variety problem", "category_id": 3, "category": {"id": 3, "name": "graph_theory", "display_name": "Graph Theory", "description": "Problems involving graphs, networks, and their properties.", "slug": "graph-theory", "order_index": 3, "created_at": "2026-07-31T15:26:25.671Z"}, "status": "open", "problem_number": "OPG-46575", "set_id": 12}
% FIRST_POSTED: 2026-10-01
% ENTRY_KIND: statement_only
% UnsolvedMath ID 3242; source catalogue code OPG-46575
% !TeX program = lualatex
% Definition and sources only; this file is not an LLM solution attempt.
\documentclass[11pt,a4paper]{article}
\usepackage[margin=25mm]{geometry}
\usepackage{fontspec}
\setmainfont{lmroman10-regular.otf}[BoldFont=lmroman10-bold.otf,
 ItalicFont=lmroman10-italic.otf,BoldItalicFont=lmroman10-bolditalic.otf]
\usepackage{amsmath,amssymb,mathtools}
\usepackage{unicode-math}
\setmathfont{latinmodern-math.otf}
\AtBeginDocument{\let\setminus\smallsetminus}
\usepackage{xurl,hyperref}
\hypersetup{colorlinks=true,urlcolor=blue}
\setcounter{secnumdepth}{0}
\setlength{\parindent}{0pt}
\setlength{\parskip}{5pt}
\setlength{\emergencystretch}{3em}
\begin{document}
\section*{Melnikov's valency-variety problem}\label{3242}
{\small UnsolvedMath ID 3242 · OPG-46575 · Graph Theory · OpenGarden}
\subsection{Definitions and mathematical statement}
Let $G=(V,E)$ be a finite simple undirected graph with $n=|V|\ge2$. Write $d(v)$ for the number of neighbors of $v$ and define its valency variety by
\[
              w(G)=|\{d(v):v\in V\}|.
\]
Thus degree values, not vertices or degree multiplicities, are counted. The chromatic number $\chi(G)$ is the least number of labels in a vertex coloring assigning distinct labels to adjacent vertices.

The proposed inequality is
\[
 \chi(G)>\left\lceil\frac{\lfloor w(G)/2\rfloor}{n-w(G)}\right\rceil
                \qquad\text{for every such }G.
\]
The floor and ceiling have their usual integer-rounding meanings. The denominator is positive: in a simple graph no isolated vertex and degree-$(n-1)$ vertex can coexist, so among the $n$ potential degrees $0,\ldots,n-1$ at least one is absent. Hence $1\le w(G)\le n-1$.

The strict inequality is essential. Since the chromatic number is integral, the assertion is equivalently a lower bound of one plus the displayed ceiling. It must hold regardless of the number of vertices realizing each degree, and no connectedness or minimum-degree condition is added. For a regular graph $w(G)=1$ the bound is only $\chi(G)>0$; the substantive regime involves many distinct degrees compared with the order. A counterexample must satisfy the exact strict comparison, rather than only a rounded or asymptotic approximation.
\subsection{Short English statement}
Does the diversity of vertex degrees force Melnikov’s stated strict lower bound on the chromatic number?
\subsection{Sources}
\begin{itemize}
\item[S1] ULAM.AI and the UnsolvedMath Contributors, supplied problems.json, record 3242 (OPG-46575), OpenGarden. Catalogue identity and source transcription; CC BY 4.0.
\url{https://huggingface.co/datasets/ulamai/UnsolvedMath}.
\item[S2] Open Problem Garden, Melnikov's valency-variety problem. Original problem and discussion.
\url{http://www.openproblemgarden.org/op/melnikovs_valency_variety_problem}.
\end{itemize}
\end{document}
